\chapter{Future Work}
\label{ch:futurework}

% DRAFT, 2026-09-16. Replaces template text from an unrelated thesis. Each item below is tied to a
% specific limitation recorded in Chapter 4 or in the project's risk register, so it should be revised
% alongside the final results rather than kept as generic future work.

\hspace{\parindent} The limitations recorded in Chapters~\ref{ch:results} and 3 point to five lines of
work that would strengthen the result or extend its reach.

\textbf{Operator control as a modelled behaviour.} The uncontrolled environment reproduces observed
bunching at the start of the corridor but accumulates about 10\% more of it over the second half.
Scheduled timepoint holding by real operators is the most plausible explanation, and it cannot be tested
here because the APC schedule field is populated on only 18\% of records. A dataset with reliable
scheduled departure times, or an operator interview protocol, would allow timepoint behaviour to be
modelled explicitly and the residual gap to be closed or attributed.

\textbf{Stop skipping as an evaluated action.} The action space includes stop skipping, and the
environment implements it with the manuscript's safeguards: never at the origin, never at the final
stop, and never at a stop the preceding bus has already skipped. Carrying riders past their stop is
counted. The reported results hold the skip action inactive so that the comparison against the two
holding rules is like for like. Evaluating the full action space, with an explicit penalty weight for
carried-past riders, is the natural next experiment.

\textbf{Weather beyond the observed support.} The estimated effect of ordinary rain on running time is
1.35\% and is not statistically significant, so the severe-weather results rest on a labelled synthetic
lognormal stress. Incident-level records, such as roadway closures or flooding reports aligned to the
corridor, would let the severe cases be grounded empirically rather than declared synthetic, and would
make the degradation curve a measurement rather than an extrapolation.

\textbf{Both directions, and other corridors.} The study models direction code 6 of Route 801. The
opposite direction is available under the same cleaning rules (226{,}233 clean stop events) and can be
used as a replication subset without any new data acquisition. Repeating the calibration and evaluation
on a second corridor would test how far the fitted parameters and the control conclusions transfer.

\textbf{Toward deployment.} All results are simulated, and no sim-to-real transfer is claimed. A
deployment study would need real-time headway and load feeds, a decision latency budget at each control
stop, and a safety review of holding instructions issued to drivers. The natural intermediate step is a
shadow evaluation in which the trained policy issues recommendations that are logged and compared with
operator decisions, without controlling any vehicle.
